\documentclass{beamer}

\mode<presentation>
{
  \usetheme{Warsaw}
  % or ...

  \setbeamercovered{transparent}
  % or whatever (possibly just delete it)
}

% \usepackage[english]{babel}
% \usepackage[latin1]{inputenc}
% \usepackage{times}
% \usepackage[T1]{fontenc}
% \usepackage{tikz}

% \usepackage{thumbpdf}
% \usepackage{hyperref}
\newcommand\yellow[1]{{\color{yellow}{#1}}} 
\newcommand\red[1]{{\color{red}{#1}}} 

% \usecolortheme{albatross}
% \setbeamertemplate{background canvas}[vertical shading][bottom=blue!100,top=blue!100]


% \setbeamercolor{normal text}{fg=white}
% \setbeamercolor{structured text}{fg=green}

\title{MOBIUS Tools Overview}
\author[J. Kiniry (UCD)]{Presented by Joseph Kiniry}



\institute{INRIA, RUN, UCD, UW\\[1ex]
  \includegraphics[width=3cm]{../../../TexInputs/mobius_transparent}}

\date{19 June 2006}

\begin{document}

\begin{frame}
  \titlepage
\end{frame}

\begin{frame}\frametitle{The MOBIUS Interactive Verification Environment}
  \begin{itemize}
  \item mandatory goals:
    \begin{itemize}
    \item develop an Eclipse-based interactive verification environment
      for reasoning about Java
    \item reason about MIDP \& JavaCard programs at the source and
      bytecode level
    \item VC generation for MOBIUS logic
    \item certificate generation
    \end{itemize}
  \item optional goals:
    \begin{itemize}
    \item integrated support for provers of interest
    \item VC generation for other logics
    \item generalized hybrid certificate generation
    \item MOBIUS component model and tool-bus
    \end{itemize}
  \item tertiary goals:
    \begin{itemize}
    \item formalized literate programming environment
    \end{itemize}
  \end{itemize}
\end{frame}

\begin{frame}\frametitle{Task 3.6 Overview}
  \begin{itemize}
  \item responsible for core MOBIUS tool analysis, design, and
    development
  \item participants: INRIA, RUN, UCD (leader), WU
  \item start date: T0+6 months
  \item end date: T0+45 months (an onward)
  \item resources: 70 man-months
  \item first deliverable: alpha version of IVE to MOBIUS partners
    at T0+12 that supports reasoning about single-threaded programs at
    the source level
    \begin{itemize}
    \item and regular component releases thereafter
    \item automatic binary and source upgrades via the Eclipse update
      manager
    \end{itemize}
  \item expect a beta version at T0+18 that supports both source and
    bytecode level reasoning about concurrent programs, certificate
    generation, and initial integrated interfaces to different proof
    assistants
  \end{itemize}
\end{frame}

\begin{frame}\frametitle{Context of Ongoing Work}
  \begin{itemize}
  \item initial component development ongoing at INRIA, RUN, UCD, and WU
  \item tight coordination with North American partners to leverage tool
    development resources
  \item independent and joint development between MOBIUS partners
  \item physical site visits for intense, focused collaboration
  \item tool and infrastructure ``working case study'' at UCD
  \item future continued virtual design and development sessions using
    VoIP, vidconf, and VNC
  \item physical collaboration in the form of site visits, tutorials
    on IVE use, component development, etc.
  \end{itemize}
\end{frame}

\begin{frame}\frametitle{The MOBIUS Process}
  \begin{itemize}
  \item a \emph{MOBIUS software development process} is being
    developed, documented and refined
    \begin{enumerate}
    \item identify an unrealized requirement(s)
    \item identify existing components that (potentially, partially)
      fulfill these requirements(s)
    \item perform a critical, detailed artifact review (measure code
      size; code complexity; documentation, specification, and unit
      test quality and coverage; design quality; etc.)
    \item decide which (sub)components to reuse for MOBIUS work
    \item update component to conform to the MOBIUS code standard
    \item add to the MOBIUS artifact repository
    \end{enumerate}
  \item once a component is part of the MOBIUS repository, evolve and
    maintain that component at a high quality level
    \begin{enumerate}
    \item automatic code standard conformance checking
    \item version control repository triggers
    \item automatic documentation generation with LP techniques
    \item PMD MOBIUS plugins (discussed later)
    \item use the MOBIUS components to develop MOBIUS components
    \end{enumerate}
  \end{itemize}
\end{frame}

\begin{frame}\frametitle{The MOBIUS Component Model}
  \begin{itemize}
  \item goals
    \begin{itemize}
    \item PCC component identification, search, retrieval
    \item dynamically compose PCC components
    \item use formal specifications to specify and reason about
      component properties
    \end{itemize}
  \item methodology
    \begin{itemize}
    \item evolving existing standard component model (from Eclipse) to
      include formal specifications and PCC
    \item use redesign and refactoring of existing tool as a case study
      in the evaluation of this new component model (RCC, discussed
      later)
    \item develop a verification component ``tool-bus'' architecture that
      incorporates PCC as a first order concept
    \end{itemize}
  \end{itemize}
\end{frame}

\begin{frame}\frametitle{Existing Components}
  \begin{itemize}
  \item JML tool suite
    \begin{itemize}
    \item type checker, compiler, runtime assertion checker, unit
      documentation generation, unit test generator
    \end{itemize}
  \item ESC/Java2
    \begin{itemize}
    \item extended static checker supporting Java 1.4 on using two
      versions of an ESC logic (unsorted $ESC_1$ and $ESC_2$) with the Simplify
      theorem prover
    \end{itemize}
  \item Jack
    \begin{itemize}
    \item type checker, VC generator for one logic and multiple provers
    \end{itemize}
  \item several provers
    \begin{itemize}
    \item Coq, PVS, Isabelle, Simplify, SMT-LIB/CVC3
    \end{itemize}
  \item American tools: Bogor, Bandara, Indus, Daikon, etc.
  \item Eclipse plugins: CodeStyle, FindBugs, PMD
  \end{itemize}
\end{frame}

\begin{frame}\frametitle{Components Under Development}
  \begin{itemize}
  \item new version of JML2 - compiler with (dramatically
    simplified) assertion semantics, improved unit test generation
  \item JML3 - type checker and compiler for Java 1.5 based upon the
    Eclipse compiler framework
  \item ESC/Java3 as a general-purpose verification framework
    \begin{itemize}
    \item specification soundness checker
    \item soundness and completeness warning system
    \item new assertion semantics for ESC
    \item logic- and prover-independent VC library
    \item new ESC logics - $ESC_2$ in PVS and Coq, $ESC_3$ in PVS and SMT-LIB
    \item VC generators for these logics
    \item (formalized, specified) prover interaction framework
    \end{itemize}
  \end{itemize}
\end{frame}

\begin{frame}\frametitle{Components Under Development (II)}
  \begin{itemize}
  \item verification-centric JML model classes
  \item ``reviving'' other SRC tools
    \begin{enumerate}
    \item Houdini - revived by RUN
    \item RCC - being revived by UCD
    \end{enumerate}
  \item BML library and VC generator
  \item Umbra bytecode library/editor
  \item Bico bytecode to logic translator
  \item INRIA Jack work integrating BML VC generator and support for
    some new provers
  \end{itemize}
\end{frame}

\begin{frame}\frametitle{Planned Components}
  \begin{itemize}
  \item (possibly multiple) VC generators for the MOBIUS logic(s)
  \item FLOSS Boogie library
  \item Jack interface extraction and reuse
  \item type checker(s) for WP2 type system(s)
  \item immutability checker
  \item correct concurrency conformance checker
  \item proof transforming compiler
  \item generalized hybrid PCC certificate library
  \item off-device checker for MOBIUS logic
  \item annotation generation framework
  \end{itemize}
\end{frame}

\begin{frame}\frametitle{Planned Components (II)}
  \begin{itemize}
  \item tertiary components
    \begin{itemize}
    \item Trusted Logic tool work (non-FLOSS)
    \item on-device checker for MOBIUS logic
    \item alternative abstract program representations (DSA, GC, graph
      representation, etc.)
    \item PMD MOBIUS plugins that summarize PCC component development
      status, quality, etc.
    \item PCC publication framework (perhaps as a Wiki plugin)
    \item formalized literate programming tool
    \item on-device checker for SMT-LIB provers and ESC logics
    \end{itemize}
  \item case studies
    \begin{itemize}
    \item MIDP applets identified in Tasks 1.3, 1.4, 3.3, etc.
    \item MOBIUS TCB components
    \item incorporation of MOBIUS PCC infrastructure into ConStruct
      pervasive framework and KOA remote voting system
    \end{itemize}
  \end{itemize}
\end{frame}

\begin{frame}\frametitle{Coordination on New Development}
  \begin{itemize}
  \item follow the MOBIUS process found (soon) in the the MOBIUS tool
    wiki closely
    \begin{enumerate}
    \item publish requirements document on MOBIUS tool wiki
    \item announce intentions on MOBIUS tool mailing list (TBC)
    \item find appropriate partners
    \item use MOBIUS IVE to develop new MOBIUS component(s)
    \item identify and encourage maintainer
    \item \emph{publish case study on tool work}
    \end{enumerate}
  \item questions or problems?
    \begin{enumerate}
    \item coordinate with UCD to ensure appropriate reuse and
      communication
    \item ask questions on the MOBIUS tool mailing list
    \item document best practices in the MOBIUS tool wiki
    \end{enumerate}
  \end{itemize}
\end{frame}

\begin{frame}\frametitle{Long-Term MOBIUS Tool Impact}
  \begin{itemize}
  \item this work has a lifespan much longer than the MOBIUS project
  \item we already have users and developers working with various
    components and many Universities using tools for instruction and
    research
  \item we eat our own dog-food to guarantee that components are
    usable and compelling
  \item we are working to ``expose'' the Science that is inherent in
    tool building of this kind to change community attitudes to
    systems-building in (applied) formal methods
  \item we are setting a high quality bar to help ensure that this
    community continues to grow and inspire others in the community
    (this is our mandate, according to Hoare)
  \end{itemize}
\end{frame}

\end{document}
